\begin{statement}{E.28}{Proved}{The complete two-adic boundary tower and its joint Ferrari lifts}
Keep the discrete groups and the ten-generator integral basis of E.26.
Put
\[
 M_H=H^{\rm ab},\qquad
 B_n=B/[H,H]H^{2^n},\qquad
 C_n=B/(K_6\cap[H,H]H^{2^n})\quad(n\ge1).
 \tag{E.28.1}
\]
Every one of the twenty-four retained $G_6$ words $f_i$ admits an
explicit correction in $N_B'$ giving compatible multiplier-one
automorphisms of every $C_n$.  One integral correcting word works
at all these levels.  The complete set of correction images in
$M_H/2^n$ coming from $N_B'$ and satisfying both hexagons has
exactly $2^n$ elements; every one is generating.  Its inverse
limit is an affine copy of $\mathbb Z_2$.

Let $\mathscr D=D_{2,6}\times_{G_6}D_6$ be the exact source
quotient of E.24.  Each of the eighteen multiplier-one points
in each of E.24's twenty-four indexed CM families extends
through the entire common tower
\[
 \mathscr C_n=\operatorname{im}
       \bigl(B\longrightarrow\mathscr D\times B_n\bigr).
 \tag{E.28.2}
\]
Thus all $432$ multiplier-one points of E.24 are retained
simultaneously with this tower.  The same single profinite
commutator word gives these compatible lifts and leaves unchanged
all six binary and all six ternary Ferrari evaluations, their
induced conjugate coordinates, both pentagons, and the generating
four-strand maps of E.24.

This tower retains commutators at every relative lower-central
depth.  If $P_n$ is the image of $N_B$ in $B_n$, then there is an
explicit primitive rank-six lattice $\mathcal D\subset M_H$ such
that
\[
 \gamma_r(P_n)=\mathcal D/2^n\mathcal D
            \simeq(\mathbb Z/2^n)^6
       \qquad(r\ge2).
 \tag{E.28.3}
\]
Consequently no quotient of $B$ with nilpotent image of $N_B$
dominates $B_n$ through the specified maps.  The assertion is
about the displayed common towers and finite Ferrari tests;
it makes no full-GT realization assertion.
\end{statement}
\begin{proof}
Index the rank-$97$ basis $n_1,\ldots,n_{97}$ as in E.26, and take
\[
 (w_1,\ldots,w_6)=
 ([n_2,n_3],[n_2,n_4],[n_2,n_5],
  [n_2,n_7],[n_2,n_9],[n_2,n_{11}]).
 \tag{E.28.4}
\]
Their integral $H$-abelianizations are the columns of
\[
 D=\begin{pmatrix}
 1&0&0&0&0&-1\\
 -2&0&0&0&1&0\\
 1&1&1&-1&-1&0\\
 1&1&0&1&-1&0\\
 0&0&0&0&1&-1\\
 -3&-1&-1&0&1&1\\
 1&-1&0&-1&0&0\\
 -1&0&0&0&0&0\\
 0&1&0&1&0&0\\
 2&0&0&1&0&-1
 \end{pmatrix}.
 \tag{E.28.5}
\]
The first six rows have determinant one.  All $4,656$ basis
commutators have integral images in $D\mathbb Z^6$.  Conjugation
by $N_B$ on $M_H$ factors through its cyclic-three deck image,
which is generated by the image of $t=n_2^{-1}$.  On these six
columns its matrix is
\[
 Q_D=\begin{pmatrix}
 0&0&0&0&-1&1\\
 0&0&1&-1&-1&1\\
 0&-2&-2&1&1&0\\
 -1&-1&-1&0&1&0\\
 1&-1&0&-1&-2&2\\
 0&-1&0&-1&-1&1
 \end{pmatrix}.
 \tag{E.28.6}
\]
The integral identities are $A_tD=DQ_D$ and
$Q_D^2+Q_D+1=0$.  Thus $D\mathbb Z^6$ is stable under every
$N_B$ conjugation.  The normal closure of all the basis
commutators is $N_B'$, so these calculations prove the exact
image equality
$\operatorname{im}(N_B'\to M_H)=\mathcal D=D\mathbb Z^6$.
The unit minor proves that reduction of $\mathcal D$ modulo
every integer embeds in the corresponding reduction of $M_H$.

Let $\mathsf S,\mathsf R$ be the integral matrices obtained by
the literal substitutions
$s:(x,y)\mapsto(y,x)$ and
$r:(x,y)\mapsto(y,(xy)^{-1})$ of E.26, and set
\[
 \mathcal E=\begin{pmatrix}1+\mathsf S\\
                         1+\mathsf R+\mathsf R^2\end{pmatrix},
 \qquad \mathcal M=\mathcal E D.
 \tag{E.28.7}
\]
For $f\in H$, the formulas
$s\mapsto sf$, $r\mapsto f^{-1}r$ satisfy the two order
relations in $B/[H,H]H^m$ exactly when
$\mathcal E[f]=0$ modulo $m$.  This follows by multiplying
$(sf)^2$ and $(f^{-1}r)^3$ in the abelian normal kernel;
their coordinates are respectively
$(1+\mathsf S)[f]$ and
$-(1+\mathsf R+\mathsf R^2)[f]$.
Their product fixes $sr$, so this is the same multiplier-one
cusp gauge as E.1.

The integral Smith factors and kernel are
\[
 \operatorname{Smith}(\mathcal M)=(1,1,1,1,3,0),\qquad
 \ker_{\mathbb Z}\mathcal M=\mathbb Z v,\quad
 v=(1,-2,1,1,1,2)^t .
 \tag{E.28.8}
\]
For every $z_i=[f_i]$ the certificate supplies an integral
$a_i\in\mathbb Z^6$ satisfying
$\mathcal M a_i=-\mathcal E z_i$.
The literal word $f_i\prod_j w_j^{(a_i)_j}$ consequently
satisfies both relations already in $B/[H,H]$.
See \texttt{bc\_cm\_fermat\_boundary\_integral.json}
for all integer exponents and resulting words.
The Smith calculation proves that modulo $2^n$ all solutions
are exactly $a_i+(\mathbb Z/2^n)v$.

Here is the actual action and the generation check.  Let $A_q$
be the deck action on $M_H$ for $q\in(\mathbb Z/3)^2$.
For the $j$th Schreier word $\eta_j$, follow its successive
deck vertices $q\to q'$ and put
\[
 \mathcal T_j=\sum_{\text{letters }y^{\pm1}\text{ in }\eta_j}
                         (A_{q'}-A_q).
 \qquad
 \mathcal A(z)=1+
       \bigl(\mathcal T_1z\ \cdots\ \mathcal T_{10}z\bigr).
 \tag{E.28.9}
\]
Once the two order relations hold, the induced substitutions
on $F$ are $x\mapsto x$ and $y\mapsto f^{-1}yf$.
Inserting the two copies of $f$ at each occurrence of
$y^{\pm1}$ gives exactly (E.28.9), including the negative
letters.  The certificate independently substitutes the
literal corrected words into every $\eta_j$, obtaining the
same integral matrix.  It also checks
$\mathcal T_jDv=0$ for every $j$.
Thus the matrix is constant on each complete two-adic
correction fibre.  Its determinants, in the saved order, are
\[
 \begin{gathered}
 1,\ 3241,\ 9937,\ 8929,\ 22681,\ 26641,\ 26353,\ 17353,\\
 17497,\ 13969,\ 25057,\ 30529,\ 447337,\ 249697,\ 408601,\ 381889,\\
 12241,\ 4249,\ 6697,\ 29017,\ 8209,\ 3097,\ 16417,\ 5977 .
 \end{gathered}
 \tag{E.28.10}
\]
All are $1$ modulo six.  In particular the matrices are
invertible modulo every $2^n$ and the action on $B/H$ is
the identity.  The induced map on $B_n$ is an automorphism.
Its $G_6$ map is unchanged because each $w_j\in N_B'$.
The resulting endomorphism of the actual common quotient
$C_n$ is injective: an element in its kernel projects to
the kernel of both component automorphisms, hence is
identity in their common image.  Finiteness gives
surjectivity.  This also proves compatibility, using the
same word at every level.

For precision at other moduli, the same restricted correction
problem has a generating solution modulo $m$ exactly when
$m$ is relatively prime to the corresponding determinant
in (E.28.10).  Sufficiency uses the displayed integral word.
For necessity, a prime divisor of that determinant is
different from $2,3$; over its residue field (E.28.8) makes
every correction differ by a multiple of $v$, which leaves
(E.28.9) unchanged and singular.  This criterion concerns
the fixed word $f_i$ and $N_B'$ corrections.  It is not an
exclusion for a whole CM matrix fibre.

We now give the exact additional comparison required for
the joint Ferrari assertion.  The map
$M_H\to F^{\rm ab}$ is $3L$, where
\[
 L=\begin{pmatrix}
 0&0&0&0&0&0&1&1&1&0\\
 0&0&1&0&0&1&0&0&0&1
 \end{pmatrix}.
 \tag{E.28.11}
\]
An integral basis matrix $Z$ for $\ker L$ has columns
\[
 e_1,e_2,e_4,e_5,e_6-e_3,e_8-e_7,e_9-e_7,e_{10}-e_3.
 \tag{E.28.12}
\]
The coordinates in this basis are the rows
$1,2,4,5,6,8,9,10$ of a vector in $\ker L$.
The following exact matrix identity supplies the correction,
without discarding the quotient by $\mathcal D$:
\[
 \mathcal M P=\mathcal E Z,\qquad
 3P=\begin{pmatrix}
 1&0&1&1&0&-1&0&0\\
 -5&-3&1&4&0&-4&0&-3\\
 2&5&-4&-6&-3&3&2&3\\
 2&2&-1&-3&0&3&2&3\\
 0&1&0&2&0&-2&1&0\\
 1&1&1&0&0&0&1&0
 \end{pmatrix}.
 \tag{E.28.13}
\]
In particular $P$ is integral over $\mathbb Z_2$.
For every actual $f\in F'$ write $[f]=Zu$.  Then
$[f]-DPu$ satisfies both boundary norms at every two-power
level.  This proves the needed compatibility for all the
Ferrari words, including their sixth-power corrections.

Generation is retained by the following literal comparison.
E.24's origins are the E.17 replacements in the indexed
CM families; they are not identified with the original
$G_6$ words $f_i$.  Reconstruct all $432$ E.24 words by
the actual formula (E.24.7), using
\texttt{de\_ferrari\_two\_family\_words.py}.
For every resulting word $f$ in row $i$, integral
Schreier rewriting verifies
\[
                 [f]-z_i\in\mathcal D+2M_H.
 \tag{E.28.13a}
\]
The explicit coordinates in the six columns (E.28.5)
are retained for every word in
\texttt{bc\_twoadic\_ferrari\_comparison.json}.
This is a congruence in the displayed boundary module,
with no equality of the two $G_6$ points assumed.
After the norm correction, the modulo-two image is
therefore one of the enumerated solutions in
$z_i+\mathcal D/2$.  Every such solution has odd
determinant.  The comparison program also checks
the rank-ten action directly for each of the $432$
corrected coordinates.  Hence the resulting action
is invertible at every two-power level.

Within a fixed E.24 source fibre, an element of $K_6$
automatically has boundary image in
$\mathcal D+2M_H$: its derived part has image
$\mathcal D$, while $n^3\in H$ and
$[n^6]=2[n^3]$ for every $n\in N_B$.
This observation is consistent with the literal
comparison, but is not used to identify the
earlier and later $G_6$ origins.

It remains to realize this correction while keeping the
entire Ferrari source fixed.  In the complete induced
binary source, the images of
$\gamma_2(N_B),\gamma_3(N_B)$ have ranks $6,0$.
This is checked by the actual conjugations of all $97$
generators on the full normal commutator image in
\texttt{de\_ferrari\_two\_lowercentral.py}.
The ternary ranks of E.23 are $15,6,3,0$ in degrees
$2,3,4,5$.  The $G_6$ image is abelian.
Consequently the actual map $B\to\mathscr D$ kills
$\gamma_5(N_B)$, in every induced coordinate.
On the boundary side,
\[
 \operatorname{Smith}(Q_D-1)=(1,1,1,3,3,3).
 \tag{E.28.14}
\]
For the six explicit words
\[
       w_j^{(5)}=[t,[t,[t,w_j]]]\in\gamma_5(N_B)
 \tag{E.28.15}
\]
their boundary columns are exactly $D(Q_D-1)^3$; the
coordinate determinant is $3^9=19683$.
This is invertible over $\mathbb Z_2$.
Hence the actual restriction of the boundary homomorphism
to $\overline{\gamma_5(N_B)}$ surjects onto
$\mathcal D\otimes\mathbb Z_2$.
Explicit preimages are products of the six words
(E.28.15) with continuous powers: embed each two-adic
exponent in $\widehat{\mathbb Z}$ with all other prime
coordinates zero.  Their images in $\mathscr D$ stay
identity.  This gives one profinite $k\in
\overline{\gamma_5(N_B)}\subset\widehat F'$ with boundary
coordinate $-DPu$, and $fk$ is the required single
commutator word for the complete tower.

We prove joint generation directly from the $B_n$
automorphism and E.24's surjective shadow
$B\to\mathscr D$.  Fix $n$, write
$W=\mathscr C_n$ and $W_N=\operatorname{im}(N_B\to W)$,
and denote the corrected homomorphism by $\psi:B\to W$.
The elements $\psi(s),\psi(r)$ belong to this actual common
image, since they are evaluated from the same profinite
word; the two order relations give the homomorphism.
Its map to $Q=B/N_B$ is the retained $G_6$ automorphism,
so $\psi(N_B)\subset W_N$.  Its $\mathscr D$ projection
is all $N_{\mathscr D}=\operatorname{im}(N_B\to\mathscr D)$:
E.24 is surjective, and an element mapping into
$N_{\mathscr D}$ has a preimage in $N_B$ by this
$Q$ comparison.  Its boundary projection is
$\mathcal A_\psi(P_n)\subset P_n$, where
$\mathcal A_\psi$ denotes the proved automorphism of $B_n$.
The two subgroups have the same finite order, so this
projection is all $P_n$.

The group $N_{\mathscr D}$ has class at most four.
Surjectivity onto $P_n$ and (E.28.3) therefore give the
exact equality
\[
 \gamma_5\bigl(\psi(N_B)\bigr)
           =\{1\}\times\mathcal D/2^n\mathcal D
           \ \subset W_N.
 \tag{E.28.16}
\]
Indeed its first projection is identity, and its second
projection is all of $\gamma_5(P_n)$; a subgroup with
trivial first coordinate is determined by its second
coordinate.  Quotienting by this subgroup, $W_N$
embeds in
$N_{\mathscr D}\times(P_n/P_n')$.
It is consequently a finite nilpotent group of primes
two and three.  It is generated by the $97$ images of
the chosen $N_B$ generators and surjects onto
$N_B^{\rm ab}/6$ through the $\mathscr D\to G_6$ map.
For each prime $p=2,3$, its Sylow Frattini quotient
thus has dimension at most $97$ and at least $97$;
the displayed projection identifies it with
$N_B^{\rm ab}/p$.  The retained $G_6$ automorphism makes
the image of $\psi(N_B)$ surjective on both Frattini
quotients.  The finite nilpotent Frattini criterion
then makes it all of the quotient of $W_N$.
Together with (E.28.16) this proves
$\psi(N_B)=W_N$.  Its surjective $Q$ map finally gives
$\psi(B)=W$.

Thus the constructed maps are compatible generating
shadows on the actual common quotients.  The definition
of $\mathscr D$ retains the six evaluations and all their
conjugate coordinates, so the correction leaves the
complete four-strand equations and generating maps
exactly as in E.24.  This proves all the joint assertions
without requiring invariance of the fixed Ferrari
source kernel.

Finally, the exact derived image is
$\mathcal D/2^n$ and commutation with $t$ induces
$Q_D-1$ on it.  Formula (E.28.14) makes this map onto.
Induction therefore gives (E.28.3) as an equality of
images of the actual lower-central subgroups.
A factorization through a quotient with nilpotent
$N_B$ image would kill one of these subgroups while
its displayed image has order $2^{6n}$, a contradiction.
All integral identities, literal source words, Smith
forms and substitution matrices are recorded by
\texttt{bc\_cm\_fermat\_boundary.py} and
\texttt{bc\_cm\_fermat\_boundary\_integral.py}, with their
executed outputs and exact JSON/array certificates.
\end{proof}
